\documentclass[10pt,a4paper]{amsart}
\usepackage[margin=19mm]{geometry}
\usepackage{amsmath,amssymb,mathtools}
\usepackage[scr=rsfs,cal=euler]{mathalfa}
\usepackage{xfrac}
\usepackage[unicode,hidelinks,bookmarks=false]{hyperref}
\newcommand{\ie}{i.e.\ }
\newcommand{\cf}{cf.\ }
\newcommand{\resp}{respectively\ }
\newcommand{\Subsection}{\subsection}
\providecommand{\nobreakdash}{\nobreak\hbox{-}}
\setlength{\emergencystretch}{2em}
\multlinegap0pt
\usepackage{amssymb}
\usepackage{mathtools}
\usepackage{xcolor}
\usepackage{bm}
\newtheorem{theorem}{Theorem}
\title{A Queueing-Theoretic Framework for Stability Analysis of LLM Inference with KV Cache Memory Constraints}
\author{Chengyi Nie, Nian Si, Zijie Zhou}
\date{Source: arXiv:2605.04595v1 (CC BY 4.0)}
\begin{document}
\maketitle

\noindent\emph{Source and licence.} A Queueing-Theoretic Framework for Stability Analysis of LLM Inference with KV Cache Memory Constraints. Version arXiv:2605.04595v1 (CC BY 4.0), distributed by arXiv under the Creative Commons Attribution 4.0 International licence, http://creativecommons.org/licenses/by/4.0/, as recorded in the arXiv licence metadata for this exact version and shown on the abstract page https://arxiv.org/abs/2605.04595v1. Full licence notice: \texttt{licenses/arxiv-2605.04595-CC-BY-4.0.txt}.\par\smallskip

\noindent\textit{Editorial scope.} Counted unit: the article's theorem thm:overload (source0.tex lines 767--769). The article's complete original proof of it is delivered with it (source0.tex lines 1103--1164, one \texttt{proof} environment of 62 lines, which the article places away from the statement). No mathematics is written, repaired or re-derived: the statement and the proof are the article's own lines, in the article's own notation, and no retained line is substituted or re-flowed.  No further source-local context is needed: every cross-reference of every retained line resolves inside the two delivered spans.  Nothing was omitted from any retained span, so the package carries no omission record.  No retained line contains a citation, so the wrapper carries no bibliography and no bibliography entry.  No retained line contains a figure environment, an image-inclusion command or drawing code, so no image is delivered and the package declares no asset dependency.  Editorial formatting only, in addition: an AMS \texttt{amsart} preamble that restates the article's own declarations and macros used by the retained lines, namely the environments \texttt{theorem} on separate counters, together with the standard packages those lines need.  The retained environments print in this document as Theorem 1 in order of appearance, whereas the article prints the same items with the numbering of its own full document.  The complete native source of the article as distributed in the arXiv e-print for this exact version is bundled unchanged as \texttt{sources/199871/source0.tex}.
\begin{theorem} \label{thm:overload}
    If $\lambda > \mu$, then under any scheduling policy, the total number of requests in the system grows to infinity. Namely, the system is overloaded.
\end{theorem}

\begin{proof}[Proof of Theorems \ref{thm:overload}]
We will show that the total outstanding memory needed in the GPU(including the chunks in the prompt phase) diverges to infinity.

For each arriving request with an input size \(s\), output length \(o\), and the chunk size \(\hat{s}\), the cumulative memory usage in the GPU (summed over the entire service of that request) is
\[
g(s,o)=(\hat{s}+2\hat{s}+\cdots+s)+((s+1)+(s+2)+\cdots+(s+o)) = \frac{(1+s/\hat{s})s+2os+(1+o)o}{2}.
\]

As the p.m.f. of joint distribution of arrival is $p(s,o)$, than the expected cumulative memory usage is 
\[
   \mathbb{E}_{s,o\sim p(s,o)}\Big[\tfrac{(1+s/\hat{s})s+2os+(1+o)o}{2}\Big] 
\]
Given that the arrival process has mean $\lambda \bar b$ in each slot, the expected total memory demand over the time period $[0,T]$ (assuming $T$ is a multiple of $\bar b$) is
\[
   \lambda T\mathbb{E}_{s,o\sim p(s,o)}\Big[\tfrac{(1+s/\hat{s})s+2os+(1+o)o}{2}\Big]. 
\]
By the strong law of large numbers for i.i.d.\ arrivals, we have
\begin{equation}
\frac{1}{T}\sum_{i=1}^{N(T)} g(s_i,o_i)
\;\rightarrow\;
\lambda \,\mathbb{E}_{s,o\sim p(s,o)}
\Big[\tfrac{(1+s/\hat{s})s+2os+(1+o)o}{2}\Big], \text{ almost surely,}
\label{LLN}
\end{equation}
where $N(T)$ denotes the number of arrivals up to time $T$, and
$s_i$ and $o_i$ are the prompt length and decode length of the $i$-th
arrival, respectively.

Now, we turn to the service side. As each batch takes \(\overline{b}\) seconds, we can take \(\tfrac{1}{\overline{b}}\) batches per second.  
Moreover, at any instant in time, the memory usage cannot exceed \(M\). This implies that the maximum total memory supply over the time period $[0,T]$  is:
\[
   T\times M
\text{ (units of memory)} 
   \times
   \frac{1}{\overline{b}}
\text{ (batches/second)}
   \;=\;
   \frac{TM}{\,\overline{b}\,}.
\]

Therefore, if \(\lambda>\mu\), we have 
\[
\lambda T\mathbb{E}_{s,o\sim p(s,o)}\Big[\tfrac{(1+s/\hat{s})s+2os+(1+o)o}{2}\Big] - \frac{TM}{\,\overline{b}\,}\rightarrow +\infty, \text{ as }T \rightarrow +\infty.
\]
Hence, by combining with Equation (\ref{LLN}), the total number of unprocessed tokens diverges to infinity almost surely.
\iffalse
\[
  \mu
  \;\times\;
  \frac{1}{2}\,\iint (s+o)(s+o+1)\,p(s,o)\,ds\,do
  \;\;\le\;\;
  \frac{M}{\,\overline{b}\,}.
\]
This implies that 
\[
  \mu
  \;\le\;
  \frac{\,2\,M\,}{\,\overline{b}\,\iint (s+o)(s+o+1)\,p(s,o)\,ds\,do}.
\]
\fi

\end{proof}

\end{document}
